\documentclass[17pt]{beamer}

\mode<presentation>

\usetheme{moloch}

\usepackage{parskip}
\usepackage{hyperref}

\title{Git}
\subtitle{in practice}

\begin{document}

\maketitle

\begin{frame}
  \centering
  \includegraphics[height=\textheight]{git-meme}
\end{frame}

\begin{frame}{Install \texttt{git}}
  \url{https://gitforwindows.org/}

  \url{https://git-scm.com/downloads}
\end{frame}

\begin{frame}
  \centering
  \includegraphics[height=\textheight]{git-for-windows}
\end{frame}

\begin{frame}{First-time setup}
  \begin{itemize}
    \item hu u?
    \item text editor
  \end{itemize}
\end{frame}

\begin{frame}{hu u?}
  \texttt{git config --global user.name <your-name>}

  \texttt{git config --global user.email <your-email>}
\end{frame}

\begin{frame}
  \centering
  \includegraphics[height=\textheight]{vscode-meme}
\end{frame}

\begin{frame}{VSCode}
  \texttt{git config --global core.editor "code --wait"}
\end{frame}

\begin{frame}{Verify}
  \texttt{git config --list}
\end{frame}

\begin{frame}{In or out?}
  \begin{itemize}
    \item tracked
    \item untracked
  \end{itemize}
\end{frame}

\begin{frame}{Tracked file states}
  \begin{itemize}
    \item modified
    \item staged
    \item committed
  \end{itemize}
\end{frame}

\begin{frame}{Create a repo}
  In project folder:

  \texttt{git init}
\end{frame}

\begin{frame}{Cloning a repo}
  \texttt{git clone <repo-url>}
\end{frame}

\begin{frame}{Status}
  \texttt{git status}
\end{frame}

\begin{frame}{Track}
  Stage modifications:

  \texttt{git add <filename>}
\end{frame}

\begin{frame}{Commit}
  \texttt{git commit}

  \texttt{git commit -m <message>}
\end{frame}

\begin{frame}
  \centering
  \includegraphics[height=\textheight]{commit-meme}
\end{frame}

\begin{frame}{Diff}
  \texttt{git diff}

  \texttt{git diff --staged}
\end{frame}

\begin{frame}{Delete}
  \texttt{git rm}
\end{frame}

\begin{frame}{Viewing history}
  \texttt{git log}
\end{frame}

\begin{frame}{Redo commit}
  \texttt{git commit --amend}
\end{frame}

\begin{frame}{Unstage}
  \texttt{git restore --staged <filename>}
\end{frame}

\begin{frame}{Restore}
  \texttt{git restore <filename>}
\end{frame}

\begin{frame}{Remotes}
  \texttt{git remote}

  \texttt{git remote add <name> <url>}
\end{frame}

\begin{frame}{Push and pull}
  \texttt{git fetch <remote>}

  \texttt{git push <remote> <branch>}
\end{frame}

\begin{frame}{Branches}
  \texttt{git branch <name>}
\end{frame}

\begin{frame}{Switch branches}
  \texttt{git checkout <branch>}
\end{frame}

\begin{frame}{Branching workflow}
  \begin{itemize}
    \item main
    \item develop
    \item feature
    \item hotfix
  \end{itemize}
\end{frame}

\begin{frame}{Branching workflow}
  \texttt{git checkout -d <name>}

  Work...
  
  \texttt{git checkout master}

  \texttt{git merge <featurebranch>}
\end{frame}

\begin{frame}
  \centering
  \includegraphics[height=\textheight]{why-fix-conflicts}
\end{frame}

\begin{frame}{Merge conflicts}
  \texttt{git merge <branch>}

  Fix conflicts...

  \texttt{git commit}
\end{frame}

\end{document}
